\begin{derivation}[从\eqnrefs{eq:sm-coleman-weinberg-si-dp12}到\eqnrefs{eq:sm-higgs-quartic-beta-dp93}]
四次耦合的 beta 函数由有效势中 $\varphi_E^4$ 的紫外对数系数决定。对任意关于背景场的多项式 $P(\varphi_E)$，定义系数提取算符
\begin{equation*}
 [\varphi_E^4]P
 \equiv\frac1{4!}\left.\frac{\dd^4P}{\dd\varphi_E^4}\right|_{\varphi_E=0}.
\end{equation*}
一圈背景质量一般可写成
\begin{equation*}
 \mathcal M_i^2(\varphi_E)=m_{i,0}^2+\kappa_i\varphi_E^2,
\end{equation*}
于是无需取大背景极限便有精确代数恒等式
\begin{equation*}
 [\varphi_E^4]\,\mathcal M_i^4(\varphi_E)=\kappa_i^2.
\end{equation*}
正文\eqnrefs{eq:sm-coleman-weinberg-si-dp12} 对重整化能标的显式导数为
\begin{equation*}
 E_{\rm R}\frac{\partial V_1}{\partial E_{\rm R}}
 =-\frac1{32\pi^2(\hbar c)^3}
 \sum_i n_i\mathcal M_i^4(\varphi_E),
\end{equation*}
所以其四次背景场系数严格为
\begin{equation*}
 [\varphi_E^4]\left(E_{\rm R}\frac{\partial V_1}{\partial E_{\rm R}}\right)
 =-\frac1{32\pi^2(\hbar c)^3}\sum_i n_\ii\kappa_i^2.
\end{equation*}
逐类计算无量纲系数 $\sum_i n_\ii\kappa_i^2$。

先算标量涨落。由树级势的 Hessian，径向 Higgs 与三个 Goldstone 模的 $\varphi_E^2$ 系数满足
\begin{equation*}
 \kappa_h=3\lambda_H,
 \qquad
 \kappa_G=\lambda_H,
 \qquad
 n_h=1,
 \qquad
 n_G=3,
\end{equation*}
因此
\begin{equation*}
 \left.\sum_i n_\ii\kappa_i^2\right|_{\rm scalar}
 =(3\lambda_H)^2+3\lambda_H^2
 =12\lambda_H^2.
\end{equation*}

再算规范涨落。$W^\pm$ 共给六个有质量物理极化自由度，$Z$ 给三个，而
\begin{equation*}
 \kappa_W=\frac{g^2}{4},
 \qquad
 \kappa_Z=\frac{g^2+g'^2}{4}.
\end{equation*}
于是
\begin{align*}
 \left.\sum_i n_\ii\kappa_i^2\right|_{\rm gauge}
 &=6\left(\frac{g^2}{4}\right)^2
 +3\left(\frac{g^2+g'^2}{4}\right)^2\\
 &=\frac38g^4+\frac{3}{16}(g^2+g'^2)^2.
\end{align*}

最后处理全部 Yukawa 矩阵，而不是预先只留下顶夸克。对任意一个 Dirac 质量本征值 $y_a$，背景质量满足
\begin{equation*}
 \mathcal M_a^2=\frac{y_a^2}{2}\varphi_E^2,
\end{equation*}
一个 Dirac 费米子给四个自旋/反粒子 Grassmann 自由度，故 $n_a=-4N_c$。其贡献为
\begin{equation*}
 n_a\kappa_a^2
 =-4N_c\left(\frac{y_a^2}{2}\right)^2
 =-N_cy_a^4.
\end{equation*}
对所有上型夸克、下型夸克和带电轻子求和，并把质量本征值平方的和改写成味基不变量，得到
\begin{align*}
 \left.\sum_i n_\ii\kappa_i^2\right|_{\rm Yukawa}
 &=-\Tr\!\left[3(Y_uY_u^\dagger)^2
 +3(Y_dY_d^\dagger)^2
 +(Y_eY_e^\dagger)^2\right]\\
 &=-\mathcal H_Y.
\end{align*}
因此全部一圈行列式的显式尺度导数由
\begin{equation*}
 \mathcal S_4
 \equiv\sum_i n_\ii\kappa_i^2
 =12\lambda_H^2
 +\frac38g^4
 +\frac{3}{16}(g^2+g'^2)^2
 -\mathcal H_Y
\end{equation*}
控制。

仅有显式对数还不够，因为正文\eqnrefs{eq:sm-effective-potential-rg-equation-dp12} 同时包含背景 Higgs 场的反常维数。在同一 $\overline{\rm MS}$ 一圈阶，最小标准模型 Higgs 双重态满足
\begin{equation*}
 16\pi^2\gamma_\Phi
 =\mathcal T_Y-\frac94g^2-\frac34g'^2.
\end{equation*}
树级四次势为
\begin{equation*}
 V_0^{(4)}
 =\frac{\lambda_H\varphi_E^4}{4(\hbar c)^3}.
\end{equation*}
把它与一圈显式尺度导数一起代入 Callan--Symanzik 方程，并只保留 $\varphi_E^4$ 与一圈量级：
\begin{equation*}
 0
 =-\frac{\mathcal S_4}{32\pi^2}
 +\frac{\beta_{\lambda_H}}4
 -\gamma_\Phi\lambda_H.
\end{equation*}
解出四次耦合的 beta 函数，
\begin{equation*}
 \beta_{\lambda_H}
 =\frac{\mathcal S_4}{8\pi^2}
 +4\gamma_\Phi\lambda_H.
\end{equation*}
乘以 $16\pi^2$ 并代入 $\mathcal S_4$ 与 $\gamma_\Phi$：
\begin{align*}
16\pi^2\beta_{\lambda_H}
={}&24\lambda_H^2
+\frac34g^4
+\frac38(g^2+g'^2)^2
-2\mathcal H_Y\\
&+4\lambda_H\mathcal T_Y
-9g^2\lambda_H-3g'^2\lambda_H\\
={}&24\lambda_H^2
-3\lambda_H(3g^2+g'^2)
+\frac38\left[2g^4+(g^2+g'^2)^2\right]
+4\lambda_H\mathcal T_Y-2\mathcal H_Y,
\end{align*}
即正文\eqnrefs{eq:sm-higgs-quartic-beta-dp93}。Higgs 双重态是 $SU(3)_c$ 单态，因此这一圈背景质量谱中没有由 Higgs 背景直接产生的胶子质量项，$\beta_{\lambda_H}$ 在该圈阶没有显式 $g_s$ 项；强耦合通过 $Y_u,Y_d$ 的重整化群方程间接进入耦合的 RG 系统。
\end{derivation}
